\section{Schnorr Instantiation}
\label{sec:schnorr}

Let $\G=\langle g\rangle$ be a cyclic group of prime order $p$, written
multiplicatively.  A secret key is $y\getsr\Z_p$ and the public key is
$Y=g^y$.  To prove knowledge of $y$, the prover samples $r\getsr\Z_p$, sends
$A=g^r$, receives $c\in\Z_p$, and returns $z=r+cy\bmod p$.  Verification checks
\begin{equation}
  \label{eq:schnorr-verify}
  g^z \stackrel{?}{=} A\,Y^c.
\end{equation}
Given accepting transcripts $(A,c,z)$ and $(A,c',z')$ with $c\ne c'$, the
witness is
\begin{equation}
  \label{eq:schnorr-extract}
  y=(z-z')(c-c')^{-1}\bmod p.
\end{equation}
Thus special soundness is perfect.  For signatures, set
$x=\mathsf{Digest}(m)$ and include the complete, canonical public-key encoding
in \cref{eq:encoding}.  The response remains one scalar, and the transform adds
exactly $\rho$ salt bits to an ordinary Schnorr signature.

\subsection{Concrete accounting}

The salt terms are information-theoretic and easy to size independently of the
group assumption.  \Cref{tab:parameters} reports conservative upper bounds for
$q_H=2^{40}$ total hash queries and several proof-query budgets.  Values are
computed directly from
\[
  \delta_{\mathrm{salt}}+\delta_{\mathrm{pre}}
  \le \frac{q_P(q_P-1)}{2^{\rho+1}}+\frac{q_Pq_H}{2^\rho}.
\]
They are bounds, not measured failure rates.

\begin{table}[t]
  \caption{Upper bounds on salt-related bad events for $q_H=2^{40}$.}
  \label{tab:parameters}
  \centering
  \small
  \begin{tabular}{@{}rrrr@{}}
    \toprule
    Salt $\rho$ & Proof queries $q_P$ & Pre-query bound & Collision bound \\
    \midrule
    128 & $2^{20}$ & $2^{-68}$ & $<2^{-89}$ \\
    128 & $2^{32}$ & $2^{-56}$ & $<2^{-65}$ \\
    192 & $2^{32}$ & $2^{-120}$ & $<2^{-129}$ \\
    192 & $2^{48}$ & $2^{-104}$ & $<2^{-97}$ \\
    256 & $2^{64}$ & $2^{-152}$ & $<2^{-129}$ \\
    \bottomrule
  \end{tabular}
\end{table}

With a $256$-bit challenge, the guessing term $2^{-t}$ is dominated by either
line of the table.  A $128$-bit salt is adequate for many bounded deployments,
but a $192$-bit salt gives more comfortable separation when the lifetime proof
count is uncertain.  The right choice should follow the system-wide query
budget, not the security level of a single key in isolation.

\subsection{Implementation profile}

The proof assumptions translate into the following interoperable profile.

\begin{enumerate}
  \item Serialize group elements canonically and reject non-canonical or
        invalid encodings before hashing or verification.
  \item Use a fixed ASCII protocol identifier and version byte.  Do not allow
        callers to select the tag dynamically.
  \item Encode all variable-length fields with fixed-width lengths, and bind
        the public key even when it is available from surrounding context.
  \item Draw each salt from the operating system's cryptographic random-number
        generator.  Treat RNG failure as a signing failure.
  \item Reduce the hash output to $\Z_p$ with a specified unbiased method, such
        as rejection sampling or a standard hash-to-field procedure.
  \item Verify lengths, group membership, and the equation in
        \cref{eq:schnorr-verify} before returning success.
\end{enumerate}

These checks are not cosmetic.  Non-canonical group encodings, for example,
can violate injectivity even when the high-level tuple in \cref{eq:encoding}
looks unambiguous.
